\begin{frame}{왜 바둑에 이런 무기가 필요했는가: 숫자}
  19$\times$19 바둑판에서 단 \textbf{81수만} 두어도 가능한
  수순은 $19^{81} \approx 10^{104}$개:
  \vskip0.5em
  \begin{itemize}
    \item 한 판(양쪽 합쳐 280$\sim$300수, 한쪽 약 80수)은 그
      트리의 \textbf{약 80개 리프}만을 탐색
    \item \textbf{100억 판}($10^{10}$)을 끝없이 두어도 도달하는
      국면은 약 $8 \times 10^{11} \approx 10^{12}$개 --- 전체
      국면 수 $10^{170}$의 $10^{-158}$ 부분
  \end{itemize}
  \vskip0.6em
  \begin{block}{완전탐색이 성립하지 않는 규모}
    ``유망한 가지만 \textbf{표본추출로 파고든다}''는 MCTS의
    원칙(15.1절 틱택토 실습의 격차의 \textbf{극한})이 사실상
    \textbf{유일한} 길 --- 그리고 그 표본추출의 방향을
    \textbf{신경망이 가중해준 PUCT + self-play 루프}가
    \textbf{이세돌 승부}를 만들어냄.
  \end{block}
  {\footnotesize 이 절 주제(목표 기반 RL과 계획)를 더 깊이
  다루는 자료: Stanford CS234.}
\end{frame}
